\begin{center}
\LARGE
The Axiom computer algebra system applied to
                 computational category theory

\end{center}

\begin{center}
\Large
Ronald Brown
\end{center}

\noindent
(Joint work with W. Dreckmann(Bangor and Stockholm)) \\
Date:  July 18th (Thursday)\\
Time:  16:25-16:45\\

\begin{center}
\large
Abstract
\end{center}

The Axiom language is based on what are called: (Axiom) categories, domains, packages, objects.
An `Axiom category' consists essentially of a signature. The representation of objects,
implementation of operations, and expression as output form, is carried out in the domain
constructors. The advantages of the Axiom language and system are discussed, and illustrated in
terms of the code for directed graphs, free categories, and the category of finite sets. It is argued
that this type of system allows for the development of code for the interaction of examples and
abstract algebraic systems, and code which is relatively easily modified, and sufficiently general to
cope with new examples. That is, the code approximates more than is usual to the standard ways of
writing mathematics. 

